% ============================================================
\subsection{Python Patterns}
% ============================================================

Good solutions to interview problems are often determined by the choice of container and knowledge of which built-in primitives express the required operation.

% ------------------------------------------------------------
\subsubsection*{Containers and cost profiles}

Selecting a container is selecting a cost profile. That decision precedes the
algorithm and frequently dominates the runtime.

\begin{center}
\begin{tabular}{lllll}
\toprule
& \textbf{Ordered} & \textbf{Mutable} & \textbf{Lookup} & \textbf{Use} \\
\midrule
\texttt{list}  & yes & yes & $O(n)$ & ordered sequence, duplicates permitted \\
\texttt{tuple} & yes & no  & $O(n)$ & fixed record, or a dictionary key \\
\texttt{set}   & no  & yes & $O(1)$ & membership, deduplication \\
\texttt{dict}  & insertion & yes & $O(1)$ & mapping from key to value \\
\bottomrule
\end{tabular}
\end{center}

\medskip

\begin{keybox}[Membership cost]
\texttt{x in some\_list} is $O(n)$. \texttt{x in some\_set} is $O(1)$ on
average. A membership test inside a loop over $n$ elements is therefore
$O(n^2)$ against a list and $O(n)$ against a set. Converting the list to a set
once, before the loop, removes a factor of $n$. This substitution accounts for
a large share of naive-to-optimal improvements.
\end{keybox}

\paragraph{Hashability.} Set elements and dictionary keys must be hashable,
which in practice requires immutability. A \texttt{tuple} is a valid key. A
\texttt{list} is not, since mutation would invalidate its position in the hash
table. Converting a list to a tuple to form a composite key is therefore a
standard construction, as in grouping anagrams by their sorted letters.

\begin{lstlisting}[style=python, caption={Tuple as a composite key}]
seen = set()
seen.add((row, col))              # legal, tuple is hashable
# seen.add([row, col])            # TypeError: unhashable type
\end{lstlisting}

\paragraph{Dictionary access.} Counting, grouping, and seen-before problems
reduce to a dictionary. Three access forms differ in their treatment of a
missing key.

\begin{lstlisting}[style=python, caption={Dictionary access}]
d[k]                  # raises KeyError if k is absent
d.get(k, default)     # returns default instead of raising
d.setdefault(k, [])   # inserts default if absent, then returns it

for k in d:                    # iterates over keys
for k, v in d.items():         # iterates over pairs
\end{lstlisting}

% ------------------------------------------------------------
\newpage
\subsubsection*{Built-in functions}

An accumulator loop is often a built-in written out longhand. The following
cover most cases.

\begin{center}
\begin{tabular}{ll}
\toprule
\textbf{Built-in} & \textbf{Expresses} \\
\midrule
\texttt{len, sum, min, max} & size and aggregation \\
\texttt{sorted} & ordering, parameterized by \texttt{key} and \texttt{reverse} \\
\texttt{any(...)} & existential quantifier: at least one element satisfies $p$ \\
\texttt{all(...)} & universal quantifier: every element satisfies $p$ \\
\texttt{enumerate} & index paired with value \\
\texttt{zip} & parallel iteration over several sequences \\
\texttt{reversed, range} & traversal order \\
\texttt{abs, round, divmod} & numeric utilities \\
\bottomrule
\end{tabular}
\end{center}

\texttt{any} and \texttt{all} short-circuit. The expression
\texttt{any(x < 0 for x in xs)} terminates at the first negative value.

\paragraph{The \texttt{key} argument.} \texttt{sorted}, \texttt{min}, and
\texttt{max} accept a \texttt{key} function applied to each element. Ranking
is performed on the image of that function while the original element is
returned. The consequence is that a record can be ranked on one of its fields
and recovered whole.

\begin{lstlisting}[style=python, caption={Ranking on a field, returning the record}]
min(scores)                            # returns a number
min(records, key=lambda r: r[1])       # returns the full record

sorted(words, key=len)                        # rank by length
sorted(records, key=lambda r: (-r[1], r[0]))  # score desc, name asc
\end{lstlisting}

A tuple key is the analogue of a multi-column \texttt{ORDER BY}. Negating a
single numeric field reverses that field alone. \texttt{reverse=True} cannot
express this, since it inverts every field simultaneously.

% ------------------------------------------------------------
\subsubsection*{Comprehensions}

A comprehension expresses selection and projection in one expression. The form
\texttt{[f(x) for x in xs if p(x)]} corresponds to selecting $f(x)$ over the
elements for which $p(x)$ holds. Four bracket types produce four containers.

\begin{lstlisting}[style=python, caption={The four forms}]
[x*x for x in xs if x > 0]        # list
{x*x for x in xs}                 # set, deduplicates
{k: v for k, v in pairs}          # dict
(x*x for x in xs)      # generator, lazy, builds no intermediate list
\end{lstlisting}

The generator form is appropriate when the result is consumed once, for
example inside \texttt{sum(...)} or \texttt{max(...)}, since it avoids
materializing a list. Nesting beyond two levels is generally less readable
than an explicit loop.

% ------------------------------------------------------------

\newpage
\subsubsection*{Sorting}

\begin{itemize}[itemsep=-0.5ex]
  \item \texttt{sorted(xs)} returns a new list. \texttt{xs.sort()} sorts in
        place and returns \texttt{None}. The assignment \texttt{xs = xs.sort()}
        binds \texttt{None} and fails silently.
  \item Python's sort is \textbf{stable}. Elements comparing equal retain
        their original relative order. Stability licenses sorting in multiple
        passes: sorting by the secondary field first and by the primary field
        second preserves the secondary ordering within each group.
  \item Sorting a sequence of lists or tuples compares elementwise from the
        left. The result depends on field order, so an explicit \texttt{key}
        should be supplied whenever the intent is nonobvious.
  \item Sorting costs $O(n \log n)$, which is negligible at the input sizes
        these problems use.
\end{itemize}

% ------------------------------------------------------------
% ------------------------------------------------------------
\subsubsection*{The \texttt{collections} module}

Four types cover nearly every use. Three are specializations of the
dictionary that eliminate the same recurring boilerplate: the check for
whether a key already exists before writing to it. The fourth replaces the
list when insertion or removal happens at the front.

\medskip

\noindent\textbf{\texttt{Counter}: frequency.}
A \texttt{Counter} is a dictionary from element to count. It is constructed in
one pass over any iterable and answers the question ``how many times does each
value occur.'' Frequency is the most common intermediate object in string and
array problems, and constructing one by hand requires an existence check on
every element.

\begin{lstlisting}[style=python, caption={Counter}]
from collections import Counter

c = Counter("mississippi")  # Counter({'i': 4, 's': 4, 'p': 2, 'm': 1})
c['s']                         # 4
c['z']              # 0, a missing key reads as zero and does NOT raise

c.most_common(2)               # [('i', 4), ('s', 4)], top 2 by count
c.most_common()             # every pair, ordered by count descending
c.total()                      # sum of all counts
\end{lstlisting}

\vspace{-2ex}

A missing key evaluates to \texttt{0} rather than raising \texttt{KeyError},
which is what permits \texttt{c[x] += 1} to be written unconditionally. The
read does not insert the key, so the counter does not grow through inspection.

\texttt{most\_common(k)} selects the top $k$ with a heap in $O(n \log k)$.
Calling \texttt{most\_common()} with no argument sorts the whole counter at
$O(n \log n)$, so the argument should be supplied whenever $k$ is known.

Counters support multiset arithmetic, which expresses several problems
directly.

\begin{lstlisting}[style=python, caption={Counter arithmetic}]
a, b = Counter("aabbc"), Counter("abbbd")

a + b        # counts added
a - b        # counts subtracted, results <= 0 are dropped
a & b        # intersection: elementwise minimum
a | b        # union: elementwise maximum
\end{lstlisting}

\begin{keybox}[Two canonical Counter identities]
\textbf{Anagram test.} Two strings are anagrams if and only if their counters
are equal, so \texttt{Counter(s) == Counter(t)} settles the question in $O(n)$.
Sorting both strings answers the same question in $O(n \log n)$.

\medskip

\textbf{Sliding-window matching.} To locate every permutation of a pattern $p$
inside a text $t$, maintain a counter over a window of length $|p|$ and
compare it against \texttt{Counter(p)} after each shift. Incrementing the
entering character and decrementing the leaving one updates the window in
$O(1)$.
\end{keybox}

\medskip

\noindent\textbf{\texttt{defaultdict}: grouping.}
A \texttt{defaultdict} is constructed with a zero-argument factory. Accessing
a missing key calls that factory, inserts the result, and returns it. The
effect is to remove the initialization branch from every accumulation loop.

\begin{lstlisting}[style=python, caption={defaultdict}]
from collections import defaultdict

groups = defaultdict(list)          # missing key -> []
for name, dept in employees:
    groups[dept].append(name)       # no existence check required

counts = defaultdict(int)           # missing key -> 0
adj    = defaultdict(set)           # missing key -> set(), an adjacency list
grid   = defaultdict(lambda: defaultdict(int))   # a nested, sparse 2D table
\end{lstlisting}

The factory is any callable taking no arguments. \texttt{int()} returns
\texttt{0}, \texttt{list()} returns \texttt{[]}, and \texttt{set()} returns an
empty set, which is why the bare type name is passed rather than a call to it.
A nested default requires a \texttt{lambda}, since
\texttt{defaultdict(defaultdict(int))} would pass an instance where a factory
is expected.

\texttt{defaultdict(list)} performs a \texttt{GROUP BY} that omits the
aggregation step. Every member of each group survives, which is the property
that \texttt{ROW\_NUMBER() OVER (PARTITION BY ...)} supplies in SQL. Building a
graph's adjacency list is the same operation applied to edges.

\begin{keybox}[Reading a \texttt{defaultdict} mutates it]
Any access to a missing key inserts it. The expression \texttt{if d[k]:}
therefore creates \texttt{k} with an empty default and the dictionary grows
through inspection alone. This corrupts a subsequent iteration over
\texttt{d.items()} and inflates \texttt{len(d)}.

For a read that must not write, use \texttt{k in d} or \texttt{d.get(k)}.
Neither triggers the factory.
\end{keybox}

The alternative to a \texttt{defaultdict} is
\texttt{d.setdefault(k, []).append(x)} on an ordinary dictionary. It produces
the same result and constructs a fresh empty list on every call, including the
calls where the key is already present. The \texttt{defaultdict} factory runs
only on a miss.

\medskip

\noindent\textbf{\texttt{deque}: a queue.}
A \texttt{deque} is a double-ended queue backed by a doubly linked list of
fixed-size blocks. Appending and popping at either end is $O(1)$. The relevant
comparison is against a \texttt{list}, which stores its elements contiguously,
so removing the front element requires shifting every remaining element left by
one position.

\begin{center}
\begin{tabular}{lll}
\toprule
\textbf{Operation} & \textbf{\texttt{list}} & \textbf{\texttt{deque}} \\
\midrule
Append or pop at the right & $O(1)$ & $O(1)$ \\
Append or pop at the left  & $O(n)$ & $O(1)$ \\
Index at position $i$      & $O(1)$ & $O(n)$ \\
\bottomrule
\end{tabular}
\end{center}

The costs are complementary. A \texttt{list} is the correct stack, since a
stack only touches the right end. A \texttt{deque} is the correct queue.
Random access into a \texttt{deque} is a sign that the wrong container was
chosen.

\begin{lstlisting}[style=python, caption={deque}]
from collections import deque

q = deque([1, 2, 3])
q.append(4)        # right
q.appendleft(0)    # left
q.pop()            # from the right
q.popleft()        # from the left, O(1) where list.pop(0) is O(n)
q.rotate(1)        # shift every element right by one, cyclically
\end{lstlisting}

Breadth-first search is the standard application. The frontier is processed in
insertion order, which requires removal from the front on every iteration of
the main loop.

\begin{lstlisting}[style=python, caption={BFS frontier}]
q = deque([start])
seen = {start}
while q:
    node = q.popleft()        # FIFO, so nodes are visited by distance
    for nxt in adj[node]:
        if nxt not in seen:
            seen.add(nxt)
            q.append(nxt)
\end{lstlisting}

\paragraph{Bounded length.} A \texttt{deque} constructed with \texttt{maxlen=k}
discards an element from the opposite end whenever a full deque is appended to.
This maintains a rolling window over the last $k$ items with no explicit
eviction logic, which is useful for moving averages and for any statistic over
a fixed lookback.

\begin{lstlisting}[style=python, caption={A rolling window of the last 3 values}]
w = deque(maxlen=3)
for x in stream:
    w.append(x)          # the oldest element is evicted automatically
    if len(w) == 3:
        print(sum(w) / 3)
\end{lstlisting}

\medskip

\noindent\textbf{\texttt{namedtuple}: a labeled record.}
A tuple carries no information about what its positions mean, so \texttt{r[1]}
requires the reader to remember the field order. A \texttt{namedtuple} attaches
names while remaining a tuple, so it stays immutable, hashable, unpackable, and
usable as a dictionary key or a set element.

\newpage

\begin{lstlisting}[style=python, caption={namedtuple}]
from collections import namedtuple

Student = namedtuple("Student", ["name", "score"])
s = Student("Ada", 92.0)

s.score            # 92.0, and s[1] still works
name, score = s    # unpacking is unchanged
min(students, key=lambda s: s.score)  # the key function reads clearly
\end{lstlisting}

The cost over a plain tuple is negligible and the gain is that a \texttt{key}
function or a sort becomes self-documenting. When mutability is required,
\texttt{dataclasses.dataclass} provides the equivalent with assignable fields.

% ------------------------------------------------------------
\subsubsection*{Iteration}

The following forms eliminate manual index arithmetic, which is the usual
source of off-by-one errors.

\begin{lstlisting}[style=python, caption={Iteration forms}]
for i, x in enumerate(xs):        # index and value together
for a, b in zip(xs, ys):          # parallel, halts at the shorter sequence
for a, b in zip(xs, xs[1:]):      # consecutive pairs
for x in reversed(xs):            # reverse traversal without index arithmetic

name, score = record              # tuple unpacking
first, *rest = xs                 # starred unpacking
\end{lstlisting}

\vspace{-2ex}

The construction \texttt{zip(xs, xs[1:])} yields every adjacent pair and
underlies run detection, difference sequences, and monotonicity checks. It is
the analogue of \texttt{LAG} in SQL.

% ------------------------------------------------------------
\subsubsection*{Problem structure}

\begin{keybox}[The four-move skeleton]
A typical warm-up problem decomposes into four steps.
\begin{enumerate}
  \item \textbf{Parse} the input into a list of records.
  \item \textbf{Reduce} the records to a target key: a minimum, a maximum, a
        count, or the $k$-th distinct value.
  \item \textbf{Filter} the records against that key.
  \item \textbf{Order and format} the output.
\end{enumerate}
Step 2 accounts for most errors. The characteristic failure is fusing steps 2
and 3 into a single pass that mishandles an edge case.
\end{keybox}

\paragraph{Ranking and ties.} The most frequent error in step 2 concerns
ties. Whether tied elements share a rank or consume one is the same decision
in Python as in SQL.

\begin{center}
\begin{tabular}{lll}
\toprule
\textbf{Tie policy} & \textbf{Python} & \textbf{SQL} \\
\midrule
Ties consume a rank & \texttt{sorted(xs)[k-1]}       & \texttt{ROW\_NUMBER() = k} \\
Ties share a rank   & \texttt{sorted(set(xs))[k-1]}  & \texttt{DENSE\_RANK() = k} \\
\bottomrule
\end{tabular}
\end{center}

For scores \texttt{[20, 20, 30]}, the second lowest is \texttt{30}. Both
elements equal to \texttt{20} hold rank 1 and neither occupies rank 2, so
\texttt{sorted(xs)[1]} evaluates to \texttt{20} and is incorrect. The
\texttt{set} converts \texttt{ROW\_NUMBER} semantics into \texttt{DENSE\_RANK}
semantics.

\paragraph{Rank then filter.} Once the key is identified, it is computed in
one pass and the records are filtered against it in a second. The structure
mirrors the necessity of wrapping a window function in a CTE, since a window
function cannot appear in \texttt{WHERE}.

\begin{lstlisting}[style=python, caption={Rank then filter}]
target  = sorted({score for _, score in records})[1]   # rank
winners = sorted(name for name, score in records if score == target)
                                            # filter, then order
\end{lstlisting}

% ------------------------------------------------------------
\vspace{-3ex}

\subsubsection*{Common errors}

\begin{itemize}
  \item \texttt{sorted(set(xs))[1]} raises \texttt{IndexError} when all values
        are identical.
  \item Slicing copies. \texttt{xs[a:b]} costs $O(b-a)$, so a slice inside a
        loop makes the loop quadratic.
  \item \texttt{list.insert(0, x)} and \texttt{list.pop(0)} are $O(n)$, since
        every subsequent element shifts. A \texttt{deque} provides $O(1)$
        operations at both ends.
  \item Integer division \texttt{//} floors toward negative infinity, so
        \texttt{-7 // 2} evaluates to \texttt{-4}. \texttt{int(-7 / 2)}
        truncates toward zero.
  \item A mutable default argument (\texttt{def f(xs=[])}) is created once at
        definition time and persists across calls. The default should be
        \texttt{None}, with the container constructed inside the func.
  \item Strings are immutable. Accumulating with \texttt{s += c} in a loop is
        $O(n^2)$. Collect the pieces in a list and apply
        \texttt{''.join(parts)}.
\end{itemize}
